% !TEX root = ../main.tex
% ============================================================================
% Owner: Tawde and Qiao (evaluation), with Aday and Bao for the Run 1 numbers.
% Drafted from confhead/results/*.md, Docs/worklog.md, and the Run 1 log.
% ============================================================================
\section{Initial Results}
\label{sec:results}

No AUSE or AURG has been measured on the full data yet; Run~1 is training at the time of writing. This section reports what the cached data says about the task and the measured training estimate, and leaves the result table and figures as placeholders to be filled when Run~1 finishes.

\subsection{Cache and label statistics}

Table~\ref{tab:cachestats} summarizes the development cache built on nuScenes v1.0-mini at a $132 \times 74$ grid and the counts measured on the full trainval cache. Three observations shape the experiments. First, the task is learnable in principle: labels cover a quarter of the coarse-grid cells, and 37\% of labeled cells are positive at $\tau = 0.10$, inside the 5--70\% gate the Run~1 plan set. Second, night is harder: the median per-image AbsRel is 0.052 by day and 0.079 at night, and the only images with no usable lidar fit are four night frames. Third, the radar residual is noisy: its median is 0.25 under the radar fit and 0.28 under the lidar fit, about five times the lidar-measured depth error, and 13.5\% of mini images have no radar fit at all. The likely cause is that nuScenes radar has no elevation, so a radar point can project onto a pixel showing another surface. Whether the channel still helps is the subject of the planned radar ablation.

\begin{table}[t]
  \centering
  \small
  \setlength{\tabcolsep}{4pt}
  \begin{tabular}{@{}l r r@{}}
    \toprule
    & \textbf{mini ($\times 2$)} & \textbf{trainval ($\times 4$)} \\
    \midrule
    Camera images & 2,424 & 204,894 \\
    Cache size & 3.43\,GB & \todo{TB} \\
    Build time & 42\,min & \todo{h} \\
    Lidar points / image & 3,613 & \todo{} \\
    Radar points / image (5 sweeps) & 150 & \todo{} \\
    Images with a lidar fit & 99.8\% & \todo{} \\
    Images with a radar fit & 86.5\% & \todo{} \\
    Labeled cells & 26.2\% & 25.1\%$^{\dagger}$ \\
    Positive labels ($\tau = 0.10$) & 37.2\% & 37.6\% \\
    Median AbsRel, day / night & 0.052 / 0.079 & \todo{} \\
    Median radar residual & 0.25 & \todo{} \\
    \bottomrule
  \end{tabular}
  \caption{Statistics of the cached head inputs and labels. $^{\dagger}$Measured on the training split at the $\times 4$ grid by the Run~1 estimate. \todo{fill the trainval column from the cache summary on CARC.}}
  \label{tab:cachestats}
\end{table}

Two label pitfalls were measured on mini and are open items (D8, D9): 3.2\% of CAM\_BACK lidar pixels fall on the ego vehicle's trunk, and about 0.35\% of lidar pixels are returns from behind a nearer object that bleed through. Both place labels on surfaces the camera does not see.

\subsection{Run 1: first full training run}

The Run~1 estimate on the trainval cache (2026-10-28) uses 631 training scenes (110,929 images, including 76 night and 125 rain scenes) and 69 dev scenes (3,000 of 12,150 images). Of the training cells, 25.1\% carry a label and 37.6\% of those are positive, giving a class weight of 1.658. The cache loads at 28 images/s while a training step runs at 1,007 images/s, so data loading sets the pace: one epoch takes about 66 minutes and 8 epochs about 8.8 hours, inside the 12-hour job limit. The run started on 2026-10-28.

\todo{Per-epoch train and dev loss, and dev AUSE per epoch, once the job ends.} On mini, the two checkpoint rules disagreed: dev AUSE was lowest after epoch 1 while dev loss was lowest after epoch 8, which is why D14 computes AUSE on the loss-selected checkpoint and reports both curves.

\begin{figure}[t]
  \centering
  \fbox{\parbox{0.9\columnwidth}{\centering\vspace{2.5cm}\todo{Run 1 loss curve: train and dev masked BCE per epoch, with dev AUSE on a second axis}\vspace{2.5cm}}}
  \caption{Run 1 training curves. \todo{replace with the figure from the run folder.}}
  \label{fig:losscurve}
\end{figure}

\subsection{Sparsification results (pending)}

Table~\ref{tab:results} and Figure~\ref{fig:sparsification} are the primary deliverable and will be filled from the Run~1 checkpoint. Every method is scored on the same labeled cells of the nuScenes val split, overall and per BEVCar condition group.

\begin{table*}[t]
  \centering
  \small
  \setlength{\tabcolsep}{3pt}
  \begin{tabular}{@{}l cc cc cc cc@{}}
    \toprule
    & \multicolumn{2}{c}{\textbf{All}} & \multicolumn{2}{c}{\textbf{Day}} & \multicolumn{2}{c}{\textbf{Night}} & \multicolumn{2}{c}{\textbf{Rain}} \\
    \cmidrule(lr){2-3}\cmidrule(lr){4-5}\cmidrule(lr){6-7}\cmidrule(lr){8-9}
    \textbf{Method} & AUSE & AURG & AUSE & AURG & AUSE & AURG & AUSE & AURG \\
    \midrule
    Random removal & -- & 0 & -- & 0 & -- & 0 & -- & 0 \\
    Temporal spread & \todo{} & \todo{} & \todo{} & \todo{} & \todo{} & \todo{} & \todo{} & \todo{} \\
    Flip disagreement & \todo{} & \todo{} & \todo{} & \todo{} & \todo{} & \todo{} & \todo{} & \todo{} \\
    Head (ours) & \todo{} & \todo{} & \todo{} & \todo{} & \todo{} & \todo{} & \todo{} & \todo{} \\
    Head, no radar & \todo{} & \todo{} & \todo{} & \todo{} & \todo{} & \todo{} & \todo{} & \todo{} \\
    Oracle & 0 & -- & 0 & -- & 0 & -- & 0 & -- \\
    \bottomrule
  \end{tabular}
  \caption{AUSE (lower is better) and AURG (higher is better) on the nuScenes val split, per BEVCar condition group. \todo{fill from Run 1; add scene counts per group.}}
  \label{tab:results}
\end{table*}

\begin{figure}[t]
  \centering
  \fbox{\parbox{0.9\columnwidth}{\centering\vspace{2.5cm}\todo{Sparsification curves: AbsRel of the remaining cells against the fraction removed, for the head, temporal spread, flip disagreement, random removal, and the oracle; one panel per condition}\vspace{2.5cm}}}
  \caption{Sparsification curves. \todo{replace with the evaluation output.}}
  \label{fig:sparsification}
\end{figure}
